% --- Archon physics formalization source begin ---
% archon:physics
% archon:covers PhyXMiniProblems/problem_phyx_mini_0475.lean
% archon:source-report reports/phyx_mini/problem_phyx_mini_0475.source.json

\chapter{Physics problem phyx\_mini\_0475}
\label{ch:PhyXMiniProblems_problem_phyx_mini_0475}

\paragraph{Problem source.}
\#\# Physical scenario

A partition divides a thermally insulated box into two compartments, each of volume \$V\$ (figure). Initially, one compartment contains \$n\$ moles of an ideal gas at temperature \$T\$, and the other compartment is evacuated. We break the partition and the gas expands, filling both compartments.

\#\# Question

What is the entropy change in this free-expansion process?

\#\# Answer choices

- A: A: nR \textbackslash{}ln 3\$

- B: B: nR \textbackslash{}ln 2\$

- C: C: nR \textbackslash{}ln 1\$

- D: D: nR \textbackslash{}ln 4\$

\#\# Figure caption (auxiliary; use the image as primary evidence)

The image contains two diagrams side by side, each depicting a container divided into sections.

**Left Diagram:**
- The container is divided into two equal sections vertically.
- The bottom section is filled with a blue substance, represented by blue dots, labeled with "V" and temperature "T."
- The top section is empty and labeled "V."
- A pink horizontal line separates the top and bottom sections, indicating a boundary or partition.

**Right Diagram:**
- The container is fully filled with the blue substance, covering both previous sections (no partition).
- The entire container is labeled "2V" with temperature "T."
- The partition seen in the left diagram is now shown as a broken pink line, indicating the division is no longer present.

Both diagrams are set against a light gray background. The blue dots likely represent gas molecules, and the diagrams visually depict a scenario of expansion or mixing, possibly illustrating an ideal gas or similar concept.

\#\# Dataset metadata

Ideal Gases and Kinetic Theory; Physical Model Grounding Reasoning, Numerical Reasoning

\paragraph{Recorded answer/context.}
B: B: nR \textbackslash{}ln 2\$

\paragraph{Figure/image path.}
/root/proposal\_for\_physic/science-mango/phyx\_mini\_run/phyx\_data/test\_image/475.png

\paragraph{Formalization target.}
create a compiling Lean file with sorry bodies at `PhyXMiniProblems/problem\_phyx\_mini\_0475.lean`. The Lean declarations must preserve the physical quantities, dimensions or dimensional roles, figure labels, governing-law hypotheses, and final relation expressed by this problem.
Use Mathlib/Physlib names found through LeanExplore where available. If a physics API is missing, introduce faithful local abstractions rather than scalar placeholder aliases.

\begin{theorem}[Physics formalization target]
\label{thm:physics:phyx_mini_0475:target}
\lean{PhyXMini0475.entropyChange_eq_nR_log_two}
\uses{def:physics:phyx-mini-0475:phyxmini0475-idealgasentropyparameters, def:physics:phyx-mini-0475:phyxmini0475-entropychange, def:physics:phyx-mini-0475:phyxmini0475-freeexpansionsetup}
The assigned autoformalize agent should translate this physics problem into Lean declarations in the covered file, with theorem and lemma proof bodies written as `by sorry`.
\end{theorem}
\begin{proof}
This is an autoformalization task, not a proof task. Produce statements that can later be proved by the physics prover without weakening the physical model.
\end{proof}

% --- Archon declaration topology begin ---
\section{Declaration topology}
\begin{definition}[volume Dimension]
  \label{def:physics:phyx-mini-0475:phyxmini0475-volumedimension}
  \lean{PhyXMini0475.volumeDimension}
  The physical dimension of volume, namely length cubed
\end{definition}
\begin{proof}
  This is the declared model definition of volume Dimension.
\end{proof}
\begin{definition}[energy Dimension]
  \label{def:physics:phyx-mini-0475:phyxmini0475-energydimension}
  \lean{PhyXMini0475.energyDimension}
  The physical dimension of energy
\end{definition}
\begin{proof}
  This is the declared model definition of energy Dimension.
\end{proof}
\begin{definition}[entropy Dimension]
  \label{def:physics:phyx-mini-0475:phyxmini0475-entropydimension}
  \lean{PhyXMini0475.entropyDimension}
  \uses{def:physics:phyx-mini-0475:phyxmini0475-energydimension}
  The physical dimension of thermodynamic entropy, energy per temperature
\end{definition}
\begin{proof}
  This is the declared model definition of entropy Dimension.
\end{proof}
\begin{definition}[Volume]
  \label{def:physics:phyx-mini-0475:phyxmini0475-volume}
  \lean{PhyXMini0475.Volume}
  \uses{def:physics:phyx-mini-0475:phyxmini0475-volumedimension}
  A volume readout in one fixed, consistent system of units
\end{definition}
\begin{proof}
  This is the declared model definition of Volume.
\end{proof}
\begin{definition}[Energy]
  \label{def:physics:phyx-mini-0475:phyxmini0475-energy}
  \lean{PhyXMini0475.Energy}
  \uses{def:physics:phyx-mini-0475:phyxmini0475-energydimension}
  An internal-energy, heat, or work readout in one fixed system of units
\end{definition}
\begin{proof}
  This is the declared model definition of Energy.
\end{proof}
\begin{definition}[Thermodynamic Entropy]
  \label{def:physics:phyx-mini-0475:phyxmini0475-thermodynamicentropy}
  \lean{PhyXMini0475.ThermodynamicEntropy}
  \uses{def:physics:phyx-mini-0475:phyxmini0475-entropydimension}
  A thermodynamic-entropy readout in one fixed system of units
\end{definition}
\begin{proof}
  This is the declared model definition of Thermodynamic Entropy.
\end{proof}
\begin{definition}[Molar Gas Constant]
  \label{def:physics:phyx-mini-0475:phyxmini0475-molargasconstant}
  \lean{PhyXMini0475.MolarGasConstant}
  \uses{def:physics:phyx-mini-0475:phyxmini0475-entropydimension}
  The molar gas constant as an entropy-per-mole quantity. Since a mole count is represented by its scalar readout, this carries the entropy dimension
\end{definition}
\begin{proof}
  This is the declared model definition of Molar Gas Constant.
\end{proof}
\begin{definition}[Ideal Gas State]
  \label{def:physics:phyx-mini-0475:phyxmini0475-idealgasstate}
  \lean{PhyXMini0475.IdealGasState}
  \uses{def:physics:phyx-mini-0475:phyxmini0475-volume, def:physics:phyx-mini-0475:phyxmini0475-energy}
  The state variables used by Physlib's monophase ideal-gas entropy function. amountMoles is the numerical amount of substance measured in moles
\end{definition}
\begin{proof}
  This is the declared model definition of Ideal Gas State.
\end{proof}
\begin{definition}[Ideal Gas Entropy Parameters]
  \label{def:physics:phyx-mini-0475:phyxmini0475-idealgasentropyparameters}
  \lean{PhyXMini0475.IdealGasEntropyParameters}
  \uses{def:physics:phyx-mini-0475:phyxmini0475-volume, def:physics:phyx-mini-0475:phyxmini0475-energy, def:physics:phyx-mini-0475:phyxmini0475-thermodynamicentropy, def:physics:phyx-mini-0475:phyxmini0475-molargasconstant}
  Parameters of Physlib's monophase ideal-gas entropy model. The reference quantities fix the entropy convention and cancel from the entropy difference
\end{definition}
\begin{proof}
  This is the declared model definition of Ideal Gas Entropy Parameters.
\end{proof}
\begin{definition}[ideal Gas Entropy]
  \label{def:physics:phyx-mini-0475:phyxmini0475-idealgasentropy}
  \lean{PhyXMini0475.idealGasEntropy}
  \uses{def:physics:phyx-mini-0475:phyxmini0475-thermodynamicentropy, def:physics:phyx-mini-0475:phyxmini0475-idealgasstate, def:physics:phyx-mini-0475:phyxmini0475-idealgasentropyparameters}
  The entropy of a state, obtained from Physlib's entropy function and tagged with its physical entropy dimension
\end{definition}
\begin{proof}
  This is the declared model definition of ideal Gas Entropy.
\end{proof}
\begin{definition}[entropy Change]
  \label{def:physics:phyx-mini-0475:phyxmini0475-entropychange}
  \lean{PhyXMini0475.entropyChange}
  \uses{def:physics:phyx-mini-0475:phyxmini0475-thermodynamicentropy, def:physics:phyx-mini-0475:phyxmini0475-idealgasstate, def:physics:phyx-mini-0475:phyxmini0475-idealgasentropyparameters, def:physics:phyx-mini-0475:phyxmini0475-idealgasentropy}
  Entropy change between two ideal-gas equilibrium states
\end{definition}
\begin{proof}
  This is the declared model definition of entropy Change.
\end{proof}
\begin{definition}[Free Expansion Setup]
  \label{def:physics:phyx-mini-0475:phyxmini0475-freeexpansionsetup}
  \lean{PhyXMini0475.FreeExpansionSetup}
  \uses{def:physics:phyx-mini-0475:phyxmini0475-volume, def:physics:phyx-mini-0475:phyxmini0475-energy, def:physics:phyx-mini-0475:phyxmini0475-idealgasstate}
  Physical and figure-derived data for the free expansion. The first-law field is a governing law; insulation and expansion into vacuum set heat and work to zero. The remaining equalities transcribe the labels V, V, 2V, T, and the conserved n moles shown or implied by the figure
\end{definition}
\begin{proof}
  This is the declared model definition of Free Expansion Setup.
\end{proof}
\begin{lemma}[internal Energy conserved]
  \label{lem:physics:phyx-mini-0475:phyxmini0475-internalenergy-conserved}
  \lean{PhyXMini0475.internalEnergy_conserved}
  \uses{def:physics:phyx-mini-0475:phyxmini0475-freeexpansionsetup}
  An insulated free expansion into vacuum conserves internal energy
\end{lemma}
\begin{proof}
  The conclusion follows from the governing relations and figure readouts named in the statement of internal Energy conserved.
\end{proof}
\begin{lemma}[final Volume eq two mul initial Volume]
  \label{lem:physics:phyx-mini-0475:phyxmini0475-finalvolume-eq-two-mul-initialvolume}
  \lean{PhyXMini0475.finalVolume_eq_two_mul_initialVolume}
  \uses{def:physics:phyx-mini-0475:phyxmini0475-freeexpansionsetup}
  Breaking the partition makes the final volume twice the initial volume
\end{lemma}
\begin{proof}
  The conclusion follows from the governing relations and figure readouts named in the statement of final Volume eq two mul initial Volume.
\end{proof}
% --- Archon declaration topology end ---
% --- Archon physics formalization source end ---
